\chapter{Conclusioni}

Nel complesso, la tesi svolta ha permesso di evidenziare come le botnet IoT rappresentino, ancora oggi, una delle minacce più concrete, pervasive e attuali per la sicurezza e la stabilità delle infrastrutture di rete globali. Sebbene il codice sorgente di Mirai e la sua architettura possano essere considerati relativamente semplici, specialmente se confrontati con soluzioni malevole più moderne che implementano crittografia avanzata o tecniche di evasione complesse, la sua efficacia operativa rimane sorprendentemente alta. Questo paradosso dimostra come tecniche apparentemente elementari, come l'attacco a forza bruta su credenziali di default, possano generare impatti devastanti se applicate su larga scala e in contesti caratterizzati da una strutturale carenza di misure di sicurezza.

L’analisi condotta mette in luce una problematica sistemica: la diffusione massiva ed esponenziale di dispositivi IoT. Questi dispositivi, che spaziano dalle telecamere di sorveglianza ai router domestici fino agli elettrodomestici intelligenti, sono spesso progettati con un focus prioritario sulla funzionalità e sul rapido utilizzo, trascurando gravemente gli aspetti di sicurezza . Questa negligenza costituisce il fattore determinante nella proliferazione di questo tipo di minacce, creando un terreno fertile per la creazione di eserciti di dispositivi "zombie" pronti a essere utilizzati. 

L’ambiente di simulazione realizzato nel corso di questo progetto, pur operando inevitabilmente su una scala ridotta rispetto alle reti reali, ha consentito di riprodurre in modo controllato e osservabile le principali dinamiche operative di una botnet. 

Questo approccio ha offerto una visione concreta e specifica dei meccanismi che regolano l'intero ciclo di vita della botnet Mirai: dalla fase di scansione e infezione, alla propagazione laterale, fino al coordinamento centralizzato tramite server Command and Control. Tale approccio si è rivelato di fondamentale importanza per superare la teoria e comprendere come le scelte progettuali adottate dagli attaccanti influenzino direttamente l'efficacia dell'attacco. Inoltre, la possibilità di osservare il comportamento della botnet in un ambiente isolato ha permesso di valutare l’impatto reale di un attacco di tipo DDoS sulla disponibilità di un servizio web, quantificando il degrado delle prestazioni e l'eventuale collasso del servizio target.Un aspetto centrale emerso con forza dal lavoro riguarda il ruolo cruciale della consapevolezza e della sicurezza preventiva. 

L'analisi del caso Mirai ha dimostrato che l’adozione di misure di protezione basilari, quali la modifica sistematica delle credenziali di fabbrica, la disabilitazione dei servizi non essenziali (come Telnet) o l’aggiornamento costante del firmware, avrebbe potuto ridurre drasticamente la superficie di attacco disponibile. Ciò evidenzia come la sicurezza dei sistemi informatici non dipenda esclusivamente dalla complessità delle soluzioni tecnologiche di difesa adottate, ma sia profondamente legata a pratiche corrette di configurazione e gestione. Questa responsabilità è condivisa: da un lato i produttori, chiamati a implementare standard di sicurezza più elevati, e dall'altro gli utenti finali e gli amministratori di rete, responsabili della corretta manutenzione dei dispositivi.

Guardando al futuro, è ragionevole e necessario ipotizzare che le botnet continueranno ad evolversi. Le future versioni di malware IoT adotteranno presumibilmente tecniche sempre più sofisticate per eludere i sistemi di rilevamento (IDS/IPS) e rendere più onerosi gli interventi di mitigazione. L’introduzione di meccanismi di comunicazione cifrati per nascondere il traffico di comando, l'adozione di architetture distribuite per eliminare i singoli punti di fallimento e l'uso di tecniche di offuscamento rappresentano sfide crescenti. In questo contesto dinamico, risulta fondamentale continuare a investire nella ricerca e nello sviluppo di strumenti di analisi e prevenzione funzionali, in grado di adattarsi a scenari di minaccia in continua mutazione. Alla luce di queste considerazioni, il lavoro presentato in questa tesi può essere considerato come una base di partenza per ulteriori approfondimenti scientifici e sviluppi tecnologici. 

L'obiettivo è quello di ampliare sia l’analisi delle botnet, sia le strategie di difesa ad esse associate. In particolare, tra le principali e direzioni di sviluppo futuro si possono individuare:
\begin{itemize} 
\item \textbf{L’analisi delle varianti evolutive}: l’estensione dello studio alle principali varianti di Mirai, come Satori e Mukashi, è fondamentale per mappare l'evoluzione del codice malevolo. Questo permetterebbe di comprendere come gli sviluppatori di malware abbiano raffinato i meccanismi di infezione, ottimizzato la propagazione e implementato nuove tecniche di persistenza per sopravvivere ai riavvi dei dispositivi, superando i limiti della versione originale;

\item \textbf{Lo studio di architetture differenti}: indagare botnet caratterizzate da architetture alternative, in particolare i modelli Peer-to-Peer (P2P) decentralizzati. A differenza del modello centralizzato queste infrastrutture non presentano un singolo punto di vulnerabilità centrale, rendendo le operazioni di smantellamento estremamente complesse e richiedendo nuove strategie di mitigazione;

\item \textbf{L’analisi dello sfruttamento di vulnerabilità note}: spostare il focus dall'analisi degli attacchi di forza bruta, come quelli utilizzati da Mirai, verso botnet che sfruttano vulnerabilità software specifiche e note (CVE). Questo permetterebbe di valutare l’impatto di exploit mirati, la rapidità con cui tali vulnerabilità vengano sfruttate dopo la loro divulgazione e come la mancata applicazione degli aggiornamenti di sicurezza influisca sulla velocità di propagazione dell'infezione;

\item \textbf{La scalabilità dell’ambiente di simulazione}: incrementare significativamente il numero di bot virtualizzati e la complessità dell'architettura simulata. L'obiettivo è valutare in modo più realistico il carico computazionale sul server bersaglio al crescere della botnet e analizzare quindi scenari più vicini a quelli del mondo reale;

\item \textbf{L’integrazione di Intelligenza Artificiale per la difesa}: l’implementazione di sistemi di rilevamento e prevenzione avanzati basati sull’analisi comportamentale e su algoritmi di Machine Learning. A differenza dei sistemi tradizionali, queste tecnologie potrebbero identificare precocemente attività malevole rilevando anomalie, permettendo di isolare i dispositivi infetti e limitare la diffusione della botnet prima che possa causare danni significativi;

\item \textbf{L'approfondimento dei protocolli IoT}: investigare i vettori di attacco che sfruttano protocolli di comunicazione specifici per l'IoT, come MQTT (Message Queuing Telemetry Transport), CoAP (Constrained Application Protocol) o ZigBee. Infatti è probabile che le future generazioni di botnet cercheranno di sfruttare questi protocolli leggeri per veicolare comandi malevoli o ottenere informazioni sensibili;

\item \textbf{Lo sviluppo di metodologie di IoT Forensics}: La definizione di framework standardizzati per l'acquisizione e l'analisi forense delle evidenze digitali sui dispositivi IoT compromessi. Ricostruendo così la catena di attacco supportando le attività di ``Incident Response'' e di attribuzione delle responsabilità.
\end{itemize}